
\section{项目概述}

本章是全书综合实战案例，实现一个完整的 \textbf{智能制造知识管理系统}，
涵盖本体设计、Java/Jena查询服务、Python推理引擎的完整技术栈。

\section{系统架构}

\begin{codebox}
\codeline{智能制造知识管理系统}
\codeline{├──\ 本体层\ (Ontology\ Layer)}
\codeline{│\ \ \ ├──\ manufacturing.owl\ \ \ \ \#\ 核心制造本体（Manchester语法）}
\codeline{│\ \ \ └──\ sparql-queries.sparql\ \#\ 标准查询集}
\codeline{├──\ 查询服务\ (Query\ Service)}
\codeline{│\ \ \ ├──\ OntologyManager.java\ \#\ 本体加载与管理（Apache\ Jena）}
\codeline{│\ \ \ └──\ QueryService.java\ \ \ \ \#\ SPARQL查询服务}
\codeline{└──\ 推理引擎\ (Reasoning\ Engine)}
\codeline{\hspace*{4\ttcw}└──\ reasoner.py\ \ \ \ \ \ \ \ \ \ \#\ 基于owlready2的推理与冲突检测}
\end{codebox}

\section{技术栈}

\begin{center}
\begin{tabular}{@{}lll@{}}
\toprule
\textbf{层次} & \textbf{技术} & \textbf{版本} \\
\midrule
本体语言 & OWL Manchester Syntax & OWL 2 DL \\
Java查询 & Apache Jena & 4.10.0+ \\
Python推理 & owlready2 + Pellet & 0.46+ \\
规则引擎 & SWRL + swrlb & - \\
\bottomrule
\end{tabular}
\end{center}

\section{快速开始}

\subsection{运行Java查询服务}

Maven依赖：
\begin{codebox}
\codeline{<dependency>}
\codeline{\hspace*{2\ttcw}<groupId>org.apache.jena</groupId>}
\codeline{\hspace*{2\ttcw}<artifactId>apache-jena-libs</artifactId>}
\codeline{\hspace*{2\ttcw}<version>4.10.0</version>}
\codeline{\hspace*{2\ttcw}<type>pom</type>}
\codeline{</dependency>}
\end{codebox}

\begin{codebox}
\codeline{mvn\ compile}
\codeline{java\ -cp\ target/classes:target/dependency/*\ QueryService}
\end{codebox}

\subsection{运行Python推理引擎}
\begin{codebox}
\codeline{pip\ install\ owlready2}
\codeline{python\ src/reasoner.py}
\end{codebox}

\section{系统核心能力}

\begin{center}
\begin{tabular}{@{}lll@{}}
\toprule
\textbf{CQ} & \textbf{问题} & \textbf{实现方式} \\
\midrule
CQ1 & 哪些设备可以加工铝合金？ & SPARQL查询 \\
CQ2 & 当前哪些设备空闲？ & SPARQL + hasStatus \\
CQ3 & 任务T001存在哪些冲突？ & SWRL规则推理 \\
CQ4 & 产品P001推荐什么工艺？ & 相似性规则 \\
CQ5 & 工单W001的工序顺序？ & 传递性属性推理 \\
\bottomrule
\end{tabular}
\end{center}

\section{开发技术说明}

\subsection{Apache Jena}
\begin{itemize}
\item \texttt{OntModel}: 支持OWL推理的本体模型
\item \texttt{ModelFactory.createOntologyModel(OntModelSpec.OWL\_DL\_MEM\_RULE\_INF)}: 创建带内置推理器的模型
\item \texttt{QueryExecutionFactory.create(sparql, model)}: 执行SPARQL查询
\end{itemize}

\subsection{owlready2}
\begin{itemize}
\item \texttt{get\_ontology(path).load()}: 加载本体
\item \texttt{sync\_reasoner\_pellet()}: 运行Pellet推理器
\item \texttt{onto.search(type=SomeClass)}: 查询特定类的实例
\end{itemize}
